% Part: proof-theory
% Chapter: natural-deduction
% Section: introduction

\documentclass[../../../include/open-logic-section]{subfiles}

\begin{document}

\olfileid{pt}{nat}{int}

\olsection{پېژندنه}

فطري استنتاج د !!{proof} د سيستمونو يوه کورنۍ ده چې په هغې کښې هر
منطقي تړونکے يا کميت ټاکونکے د استنتاجونو يوه جوړه لري: يو استنتاج
تړونکے «داخلوي» (د !!a{introduction} قاعده) او بل ئې «وباسي»
(د !!a{elimination} قاعده)۔ د بېلګې په توګه، د \Intro{\lif}
پايله د~$!A \lif !B$ په بڼه !!{formula}s لري او د \Elim{\lif}
مقدمې د~$!A \lif !B$ په بڼه !!{formula}s لري۔

د فطري استنتاج لومړۍ دوه بڼې په ۱۹۳۰مو کلونو کښې ګېرهارډ
ګېنټسن او ستانيسواف ياشکووسکي وړاندې کړې۔ د ثبوت د تيورۍ پوهان
تر ډېره د ګېنټسن سيستم څېړي؛ د ياشکووسکي سيستم بيا هغو سيستمونو
ته لار هواره کړه چې په تدريس کښې تر ټولو ډېر کارېږي۔ په دواړو
سيستمونو کښې له \emph{فرضيو} څخه پيل کولے شو: داسې !!{formula}s
چې !!{prove}d کېدو ته اړتيا نۀ لري، خو له هغو څخه د نورو
\psOblique{!!{formula}s} په !!{proving} کښې کار اخيستلے شو۔ ټول !!{proof} په
داسې فرضيو پورې تړلے کېدے شي۔ دا چې فرضيې ثابتې نۀ دي، نو
!!{proof} خپله پايله (خپل \emph{پاييز !!{formula}}) نۀ ثابتوي،
بلکې يوازې دا ثابتوي چې پايله له فرضيو څخه راځي۔

د استنتاج ځينې قاعدې داسې دي چې پايلې ئې بيا په ځينو فرضيو
پورې نۀ تړل کېږي: دغه فرضيې «وباسي»۔ د بېلګې په توګه،
\Intro{\lif} د فرضيې~$!A$ په مرسته د پايلې~$!B$ يو
!!a{proof} اخلي او پايله ئې $!A \lif !B$ وي۔ په~$!A \lif !B$ ختمېدونکے
!!{proof} بيا په~$!A$ پورې نۀ تړل کېږي؛ فرضيه $!A$ ايستل
شوې ده۔ د ګېنټسن په سيستمونو کښې !!{proof}s د
\psOblique{!!{formula}s} ونې دي، فرضيې تر ټولو بره !!{formula}s
دي، او د~\Intro{\lif} غوندې قاعدې داسې نښې لري چې ايستل شوې
فرضيې ښيي۔ د ياشکووسکي په سيستمونو کښې د !!{proof} هغه برخې
چې په يوه فرضيه پورې تړلې وي، په يو ډول بېلې ښودل کېږي؛ مثلاً
له يوې ولاړې کرښې وروسته لږ دننه ليکل کېږي يا په يوه بکس کښې
راځي۔ ايستل کېدونکې فرضيه~$!A$ او د شرط تړنې تالي $!B$
په بکس کښې لومړۍ او وروستۍ کرښه وي، او د \Intro{\lif}
پايله~$!A \lif !B$ له بکس څخه بهر وي۔

دغه سيستمونه ځکه «فطري» دي چې د استنتاج قاعدې ئې زموږ
(يا لږ تر لږه د رياضي پوهانو) د فطري استدلال ډول ښيي۔ د يوه
فرضي حاصل (يوې شرط تړنې) ثابتولو دپاره مقدم فرضوو او په هغې
تالي ثابتوو۔ همدا \Intro{\lif} ده۔ چې شرط تړنه $!A \lif !B$ او
د هغې مقدم~$!A$ راته معلوم وي، پايله~$!B$ اخلو۔ همدا
\Elim{\lif} ده۔ د حالتونو له مخې ثبوت په دې کاروو چې وښيو
د فصل تړنې~$!A \lor !B$ په فرضولو يوه پايله راځي۔ همدا \Elim{\lor}
ده۔ د عمومي دعوې~$\lforall[x][!A(x)]$ ثابتولو دپاره ثابتوو چې د يوه
خپل‌سري~$c$ دپاره $!A(c)$ رښتيا ده۔ همدا \Intro{\lforall}
ده، او داسې نور۔

د بېلګې په توګه، د ګېنټسن په طرز فطري استنتاج کښې
د~$!A \lif (!A \lor !B)$ يو ساده !!{proof} داسې دے:
\[
\AxiomC{$\Discharge{!A}{x}$}
\RightLabel{\Intro{\lor}}
\UnaryInfC{$!A \lor !B$}
\DischargeRule{\Intro{\lif}}{x}
\UnaryInfC{$!A \lif (!A \lor !B)$}
\DisplayProof
\]
دا له فرضيې~$!A$ څخه پيلېږي، په \Intro{\lor} ئې
پايله~$!A
\lor !B$ اخلي، او بيا په \Intro{\lif} ئې
پايله~$!A \lif (!A \lor !B)$ اخلي۔ د \Intro{\lif} استنتاج
فرضيه~$!A$ وباسي؛ نښه~$x$ همدا ښيي۔

د ثبوت د تيورۍ پوهان د ګېنټسن اصلي سيستمونه غوره کوي چې د
\psOblique{!!{formula}s} پر ونو کار کوي، خو يو بدل ډول ئې هم
د ثبوت په تيورۍ کښې مهم دے۔ !!^a{proof} چې له فرضيو~$\Gamma$
څخه~$!A$ ثابتوي، ښيي چې $!A$ له~$\Gamma$ څخه راځي؛ يعنې
$\Gamma \Entails !A$۔ د فطري استنتاج قاعدې په فطري ډول د دې ډول منطقي
لازميتونو په هکله لوستل کېدے شي۔ مثلاً \Intro{\lif} وايي چې
که $\Gamma + !A
\Entails !B$ وي، نو $\Gamma \Entails !A \lif !B$ دے۔ دا فطري ده چې قاعدې نېغ په دې
ډول جوړې کړو چې د استنتاج د قاعدې په مقدمو او پايله کښې پر
فرضيو او له هغو پر راوتونکې پايله دواړو کار وکړي؛ يعنې پر
سېکوېنټونو $\Gamma \Sequent !A$ کار وکړي۔ پورته ساده ثبوت په داسې سيستم
کښې داسې ښکاري:
\[
\Axiom$x:!A \fCenter !A$
\RightLabel{\Intro{\lor}}
\UnaryInf$x:!A \fCenter !A \lor !B$
\DischargeRule{\Intro{\lif}}{x}
\UnaryInf$\fCenter !A \lif (!A \lor !B)$
\DisplayProof
\]
لکه څنګه چې وينئ، قاعدې هماغه دي: د~$\Sequent$ ښي لاس پر !!{formula}
کار کوي؛ چپ لاس !!{formula}s يوازې هغه فرضيې ښيي چې
هر ګام په هغو پورې تړلے دے۔
\textbf{سرچينه‌يي سمون OLCMP-134:} د شرط تړنې د داخلولو په
پايله کښې مقدم نۀ پرېوځي؛ پايله هماغه شرط تړنه ده۔

د !!{introduction}/!!{elimination} د قاعدو ټولې جوړې په
خپل سر د کلاسيکي منطق دپاره بشپړې نۀ دي۔ دوه نورې قاعدې
ورزياتول غواړي چې په~$\lfalse$ پورې اړه لري۔ لومړۍ
«د شهودي محال قاعده»~\FalseInt، يا «چاودنه» ده، چې وايي
له $\lfalse$ څخه هر څه استنتاجولے شو۔ دويمه د غيرمستقيم ثبوت
کلاسيکي اصل~\FalseCl{} (يا «د کلاسيکي محال قاعده») ده،
چې وايي که له~$\lnot !A$ څخه~$\lfalse$ !!{prove} کولے شو، نو
$!A$ مو !!{prove}d کړے دے۔ که \FalseCl پرېږدو، د شهودي
منطق مناسب سيستم ترې جوړېږي۔ که دواړه پرېږدو، هغه سيستم
ترې جوړېږي چې «لږ تر لږه منطق» نومېږي۔

موږ به د کلاسيکي او شهودي منطق دپاره د ګېنټسن په طرز
فطري استنتاج وڅېړو۔ دا د \Log{N1c} او \Log{N1i} سيستمونه
دي۔ د \Log{N2c} او \Log{N2i} سيستمونه ئې ورته بڼې دي چې
د يوازېنيو \psOblique{!!{formula}s}~$!A$ پر ځاے پر
سېکوېنټونو $\Gamma \Sequent !A$ کار کوي۔

\end{document}
